\chapter{Hodge-родитель градуировки полного пространства рёбер}
\label{ch:v7-edge-grading-hodge-superconnection-parent-gate}

\section{Вопрос родительского происхождения}

Циклически-изотипический гейт построил проектор $P_*$ ранга шесть на
пространстве одиннадцати новых стрелок и инволюцию
$\Gamma_E=I-2P_*$. Теперь необходимо получить не только желаемые знаки, но
и их квартичное насыщение из одной нормы отображения момента.

Положим
\begin{equation}
 P_+=P_*,\qquad P_-=I-P_*,\qquad
 \Gamma_E=P_--P_+,qquad
 (\rank P_+,\rank P_-)=(6,5).
 \label{eq:v7-edge-hodge-projectors}
\end{equation}
Индексы $+$ и $-$ здесь обозначают целевую и нежелательную опоры, а не
знак будущей квадратичной формы.

\section{Двухступенный комплекс пространства стрелок}

Возьмём две копии полного пространства новых стрелок
\begin{equation}
 \mathcal K_E=\mathcal E_E^0\oplus\mathcal E_E^1,
 \qquad
 \dim_{\mathbb C}(\mathcal E_E^0,\mathcal E_E^1)=(11,11),
 \qquad
 \chi_E=\operatorname{diag}(-I_{11},I_{11}).
 \label{eq:v7-edge-hodge-two-term-carrier}
\end{equation}
Для амплитуд новых рёбер $z=(z_e)$ определим диагональную стрелку
$Z=\operatorname{diag}(z_e)$ и динамический дифференциал
\begin{equation}
 d_Z=\begin{pmatrix}0&0\\Z&0\end{pmatrix},
 \qquad d_Z^2=0.
 \label{eq:v7-edge-hodge-dynamic-differential}
\end{equation}

Второй, фиксированный цвет стрелок задаётся уже выведенными проекторами:
\begin{equation}
 \delta_E=\begin{pmatrix}0&P_+\\P_-&0\end{pmatrix},
 \qquad \delta_E^2=0.
 \label{eq:v7-edge-hodge-background-differential}
\end{equation}
Нильпотентность следует из $P_+P_-=0$. Это не новое динамическое поле:
$\delta_E$ является фоновым дифференциалом, производным от точного
циклически-изотипического разбиения.

Его Hodge-момент равен самой градуировке:
\begin{equation}
 [\delta_E,\delta_E^\dagger]
 =-\widehat\Gamma_E,
 \qquad
 \widehat\Gamma_E=\operatorname{diag}(\Gamma_E,-\Gamma_E).
 \label{eq:v7-edge-hodge-derived-grading}
\end{equation}
Таким образом, сдвиг не постулируется отдельной матрицей: он является
квадратичным следствием нильпотентной фоновой стрелки.

\section{Одно отображение момента}

Рассматриваем $d_Z$ и $\delta_E$ как два ортогональных цвета стрелок одного
представления колчана. Для пространства представлений колчана отображение
момента является суммой вкладов стрелок, а квадрат его нормы задаёт
стандартный Morse-функционал \cite{v7-edge-hodge-quiver-moment}.
При общем масштабе $\mu>0$ получаем
\begin{equation}
 \mathfrak m_\mu(Z)
 =[d_Z,d_Z^\dagger]
 +\mu^2[\delta_E,\delta_E^\dagger]
 =[d_Z,d_Z^\dagger]-\mu^2\widehat\Gamma_E.
 \label{eq:v7-edge-hodge-total-moment}
\end{equation}
Единственный функционал равен
\begin{equation}
 \mathcal S_\mu(Z)
 =\frac12\Tr_{\mathcal K_E}\mathfrak m_\mu(Z)^2
 -5\mu^4.
 \label{eq:v7-edge-hodge-action}
\end{equation}
Последнее слагаемое есть не динамический вес, а вычитание постоянной
энергии пяти положительных фоновых блоков.

Прямая подстановка даёт точное разложение
\begin{equation}
 \boxed{
 \mathcal S_\mu(z)
 =\sum_{e\in E_*}(|z_e|^2-\mu^2)^2
 +\sum_{e\notin E_*}(|z_e|^4+2\mu^2|z_e|^2).}
 \label{eq:v7-edge-hodge-reduced-action}
\end{equation}
Оба типа квадратичных знаков и все квартичные коэффициенты следуют из одной
неотрицательной нормы. Независимый стабилизирующий вес отсутствует.

\section{Запуск и точный вакуум}

Нуль является стационарной точкой. Его вещественный гессиан на одном
поколении имеет сигнатуру
\begin{equation}
 (n_-,n_0,n_+)_{z=0}=(12,0,10).
 \label{eq:v7-edge-hodge-origin-hessian}
\end{equation}
Двенадцать отрицательных направлений принадлежат шести целевым комплексным
рёбрам; десять положительных --- пяти нежелательным.

Из \eqref{eq:v7-edge-hodge-reduced-action} множество глобальных минимумов
задаётся без решения связанных уравнений:
\begin{equation}
 |z_e|^2=\mu^2\quad(e\in E_*),
 \qquad
 z_e=0\quad(e\notin E_*),
 \qquad
 \min\mathcal S_\mu=0.
 \label{eq:v7-edge-hodge-vacuum}
\end{equation}
В единичном представителе вакуумный гессиан имеет сигнатуру
\begin{equation}
 (n_-,n_0,n_+)_{z=z_*}=(0,6,16).
 \label{eq:v7-edge-hodge-vacuum-hessian}
\end{equation}
Шесть нулей являются фазами выбранных рёбер; радиальные моды и все
нежелательные отклонения положительны.

\section{Семейный подъём}

Если каждое $z_e$ заменить семейной матрицей $Z_e\in M_3(\mathbb C)$, то
тот же незавешенный след даёт
\begin{equation}
 Z_e\in\mu U(3)\quad(e\in E_*),
 \qquad
 Z_e=0\quad(e\notin E_*).
 \label{eq:v7-edge-hodge-family-vacuum}
\end{equation}
Следовательно, все шесть целевых матриц имеют полный ранг, а пять
конкурирующих исчезают. Сигнатуры равны
\begin{equation}
 (n_-,n_0,n_+)_{0}^{\rm family}=(108,0,90),
 \qquad
 (n_-,n_0,n_+)_{*}^{\rm family}=(0,54,144).
 \label{eq:v7-edge-hodge-family-hessians}
\end{equation}
Оставшиеся $54$ нуля образуют касательные направления $U(3)^6$. Действие
выбирает опору и сингулярные числа, но ещё не относительные семейные
ориентации.

\section{Real-завершение и граница физического статуса}

Самосопряжённые операторы $d_Z+d_Z^\dagger$ и
$\delta_E+\delta_E^\dagger$ нечётны относительно $\chi_E$. Антилинейный
обмен двух ступеней переводит ориентированную стрелку в сопряжённую и
удовлетворяет
\begin{equation}
 J_E\widehat\Gamma_EJ_E^{-1}=-\widehat\Gamma_E.
 \label{eq:v7-edge-hodge-real-grading}
\end{equation}
Полное Real-удвоение удваивает след, а физический полуслед восстанавливает
\eqref{eq:v7-edge-hodge-action} без нового коэффициента. Такое чтение
соответствует общей суперсвязностной идее нечётного нуль-форменного
эндоморфизма и его кривизности \cite{v7-edge-hodge-quillen}.

\begin{theorem}[Полевой Hodge-родитель точного меню]
На двухступенном комплексе полного пространства стрелок фиксированный
нильпотентный дифференциал $\delta_E$ выводит градуировку $\Gamma_E$ своим
Hodge-моментом. Одна норма суммарного отображения момента
\eqref{eq:v7-edge-hodge-action} делает нуль неустойчивым ровно по шести
целевым рёбрам, даёт положительную щель пяти нежелательным и имеет
поперечно устойчивые минимумы, в которых все шесть целевых блоков ненулевые.
Квадратичный запуск и квартичное насыщение не имеют независимых весов.
\end{theorem}

Результат строг на вспомогательном пространстве полей. Ещё не доказано,
что два цвета стрелок являются каноническими компонентами конечной
Real-бимодульной суперсвязности физической алгебры. Масштаб $\mu$ не выведен,
а семейное вакуумное многообразие $U(3)^6$ не сведено к наблюдаемым
смешиваниям.

\begin{nogocriterion}[Запрет преждевременного полного замыкания]
Нельзя считать \eqref{eq:v7-edge-hodge-action} физическим спектральным
действием до вывода обоих цветов стрелок из одного Real-бимодуля, проверки
первого порядка, кинетической метрики и gauge-quotient полного гессиана.
Нельзя также выбирать шесть матриц из $U(3)^6$ по наблюдаемым CKM/PMNS.
\end{nogocriterion}

\begin{protocol}
\begin{enumerate}
 \item двухступенный полевой носитель: \textbf{$11+11$};
 \item динамический дифференциал $d_Z^2$: \textbf{$0$};
 \item фоновый дифференциал $\delta_E^2$: \textbf{$0$};
 \item градуировка из фонового момента: \textbf{получена};
 \item одна норма отображения момента: \textbf{получена};
 \item независимый квартичный вес: \textbf{нет};
 \item нулевая сигнатура на поколение: \textbf{$(12,0,10)$};
 \item вакуумная сигнатура на поколение: \textbf{$(0,6,16)$};
 \item целевых ненулевых рёбер: \textbf{$6$};
 \item нежелательных ненулевых рёбер: \textbf{$0$};
 \item семейный вакуум: \textbf{$\mu U(3)^6$};
 \item семейная вакуумная сигнатура: \textbf{$(0,54,144)$};
 \item Real-полуслед: \textbf{коэффициенты сохраняет};
 \item физическое бимодульное вложение: \textbf{открыто};
 \item итог: \textbf{положительный полевой Hodge-родитель}.
\end{enumerate}
\end{protocol}

Машинный сертификат сохранён в
\path{s2t_v7_edge_grading_hodge_superconnection_parent_gate_results.json}.

\begin{thebibliography}{9}
\bibitem{v7-edge-hodge-quiver-moment}
M.~Harada, G.~Wilkin,
\emph{Morse theory of the moment map for representations of quivers},
Geometriae Dedicata 150 (2011), 307--353; arXiv:0807.4734.
\bibitem{v7-edge-hodge-quillen}
D.~Quillen,
\emph{Superconnections and the Chern character},
Topology 24 (1985), 89--95.
\end{thebibliography}